\clearpage
\fancyhead[L]{\small GAUSS AN BODE.}
\hypertarget{native-459-bode}{}
\begin{center}{\large GAUSS AN BODE.}\end{center}

\begin{center}
\rule{0.55\linewidth}{0.4pt}\\[2pt]
Astronomisches Jahrbuch für 1829. Seite 144\,..\,145. Berlin 1826.\\[2pt]
\rule{0.55\linewidth}{0.4pt}
\end{center}

\begin{flushright}
Göttingen, den 10. Juli 1826.
\end{flushright}

\begingroup\rightskip=0pt plus 1em\relax
Hiermit sende ich Ihnen, mein hochverehrter Freund, meine Meridianbeobachtungen der Pallas und der Ceres um die Zeit ihrer diesjährigen Oppositionen zur freien Verwendung in Ihrem Jahrbuch.
\par\endgroup

\begin{center}
\emph{Beobachtungen der Pallas} 1826.
\end{center}

\begin{center}
\begingroup\small\setlength{\tabcolsep}{4pt}
\begin{tabular}{lr|c|c|c}
 & & \shortstack{Mittlere Zeit\\in Göttingen} & Rektaszension & \shortstack{Deklination\\nördl.} \\
\hline
Juni & 21 & $12^{\mathrm u}11^{\mathrm m}35^{\mathrm s}4$ & $272^\circ 31' 4''9$ & $23^\circ 52' 19''4$ \\
 & 22 & $12\ \ 6\ \ 48.7$ & $272\ \ 18\ \ 21.3$ & $23\ \ 51\ \ 24.1$ \\
 & 23 & $12\ \ 2\ \ \ 2.9$ & $272\ \ \ 5\ \ 50.8:$ & $23\ \ 50\ \ \ 1.9$ \\
 & 24 & $11\ 57\ 16.2$ & $271\ 53\ \ 7.0$ & $23\ 48\ 32.4$ \\
 & 25 & $11\ 52\ 30.0$ & $271\ 40\ 30.6$ & $23\ 46\ 39.4$ \\
 & 26 & $11\ 47\ 43.6$ & $271\ 27\ 51.1$ & $23\ 44\ 24.3$ \\
 & 27 & $11\ 42\ 58.0$ & $271\ 15\ 23.9:$ & ---\quad---\quad--- \\
 & 28 & $11\ 38\ 12.4$ & $271\ \ 2\ 55.8$ & $23\ 39\ 10.5$ \\
\end{tabular}
\endgroup
\end{center}

Das schwache Licht des Planeten erschwerte die Beobachtungen und verringerte ihre Genauigkeit. Am 27. ließ ich mich durch einen Stern 10. Größe irreführen, der nur eine Bogenminute von Pallas entfernt stand. Dadurch wurde die Zenitdistanz nicht auf den erheblich schwächeren Planeten eingestellt, und dieser wurde außerdem nur an zwei Fäden unzuverlässig beobachtet.

\Needspace{15\baselineskip}
\SourceLeaf{460}

\begin{center}
\emph{Beobachtungen der Ceres} 1826.
\end{center}

\begin{center}
\begingroup\small\setlength{\tabcolsep}{4pt}
\begin{tabular}{lr|c|c|c}
 & & \shortstack{Mittlere Zeit\\in Göttingen} & Rektaszension & \shortstack{Deklination\\südl.} \\
\hline
Juni & 24 & $12^{\mathrm u}30^{\mathrm m}16^{\mathrm s}6$ & $280^\circ 9' 33''0:$ & $27^\circ 40' 54''6$ \\
 & 25 & $12\ 25\ 22.6$ & $279\ 54\ 59.5$ & $27\ 45\ 11.6$ \\
 & 26 & $12\ 20\ 29.0$ & $279\ 40\ 32.7$ & $27\ 49\ 20.4$ \\
 & 27 & $12\ 15\ 35.2$ & $279\ 26\ \ 1.5$ & $27\ 53\ 25.9$ \\
 & 28 & $12\ 10\ 41.1$ & $279\ 11\ 26.4$ & $27\ 57\ 26.8$ \\
 & 30 & $12\ \ 0\ 52.3$ & $278\ 42\ \ 7.6$ & $28\ \ 5\ 20.8$ \\
Juli & \ 2 & $11\ 51\ \ 3.6$ & $278\ 12\ 49.0$ & $28\ 12\ 57.7$ \\
 & \ 3 & $11\ 46\ \ 9.3$ & $277\ 58\ 11.1$ & $28\ 16\ 37.1$ \\
\end{tabular}
\endgroup
\end{center}

Ceres hatte gut die 8. Größe, und alle Beobachtungen sind gut. Nur die erste Rektaszension beruht auf zwei Fäden allein, und dabei wurde die schwächste Vergrößerung verwendet.
